% Preamble of the course "Mowing a real lawn with a car-like robot"
\usepackage[a4paper,margin=2.5cm]{geometry}
\usepackage[T1]{fontenc}
\DeclareUnicodeCharacter{00B0}{\ensuremath{^\circ}}
\usepackage{lmodern}
\usepackage{microtype}
\usepackage{amsmath,amssymb,mathtools}
\usepackage{graphicx}
\usepackage{booktabs}
\usepackage{enumitem}
\usepackage{siunitx}
\sisetup{group-separator={\,}, detect-all}
\usepackage[dvipsnames]{xcolor}
\usepackage{tikz}
\usetikzlibrary{calc,arrows.meta,angles,quotes,decorations.pathreplacing,
  positioning}
\usepackage[most]{tcolorbox}
\usepackage[colorlinks,linkcolor=blue!50!black,citecolor=blue!50!black,
  urlcolor=blue!50!black]{hyperref}

\setlist{nosep,leftmargin=*}
\setlength{\parindent}{0pt}
\setlength{\parskip}{0.6em}

% ---------------------------------------------------------------- boxes
% blue: definitions, red: theorems, amber: propositions, teal: "In words",
% grey: examples and remarks
\newtcbtheorem[number within=chapter]{definition}{Definition}{
  enhanced, breakable, colback=blue!4, colframe=blue!55!black,
  fonttitle=\bfseries, separator sign={.}}{def}
\newtcbtheorem[use counter from=definition]{theorem}{Theorem}{
  enhanced, breakable, colback=red!4, colframe=red!60!black,
  fonttitle=\bfseries, separator sign={.}}{thm}
\newtcbtheorem[use counter from=definition]{proposition}{Proposition}{
  enhanced, breakable, colback=orange!6, colframe=orange!75!black,
  fonttitle=\bfseries, separator sign={.}}{prop}
\newtcbtheorem[use counter from=definition]{example}{Example}{
  enhanced, breakable, colback=black!3, colframe=black!45,
  fonttitle=\bfseries, separator sign={.}}{ex}
\newtcolorbox{inwords}{enhanced, breakable, colback=teal!5,
  colframe=teal!60!black, fonttitle=\bfseries, title={In words}}
\newtcolorbox{incode}{enhanced, breakable, colback=black!2,
  colframe=black!30, fonttitle=\bfseries, title={In the code}}

\newenvironment{proof}{\par\textit{Proof.}\ }{\hfill$\square$\par}

% ------------------------------------------------------------- notation
\newcommand{\R}{\mathbb{R}}
\newcommand{\kmax}{\kappa_{\max}}
\newcommand{\Rmin}{R_{\min}}
\newcommand{\weff}{w_{\mathrm{eff}}}
\newcommand{\code}[1]{\texttt{#1}}
\DeclareMathOperator{\atantwo}{atan2}
